\documentclass[11pt,reqno]{amsart}
\usepackage[T1]{fontenc}
\usepackage{amsmath,amssymb,amsthm,mathtools,booktabs,array}
\usepackage[margin=1in]{geometry}
\usepackage[colorlinks=true,linkcolor=blue,urlcolor=blue]{hyperref}
\newtheorem{proposition}{Proposition}[section]
\newtheorem{theorem}[proposition]{Theorem}
\theoremstyle{remark}\newtheorem{remark}[proposition]{Remark}
\newcommand{\R}{\mathbb R}\newcommand{\F}{\mathcal F}
\newcommand{\Id}{I}\newcommand{\NN}{\mathcal N}
\DeclareMathOperator{\diag}{diag}
\title[Finite Spin, Krawtchouk Systems, and Oscillators]{From Finite Spin to the Harmonic Oscillator:\\Krawtchouk Transforms, Mehler Kernels, and Convolution}
\date{A synthesis of FeinNotes11, the Kravchuk Encyclopedia,\\the accompanying research notes, and the present discussion}
\begin{document}
\begin{abstract}
We give a convention-explicit account of the passage from finite spin to
Hermite functions and the harmonic oscillator. Three structures are kept
separate: the finite spin representation of $\mathfrak{sl}_2$, its
Heisenberg contraction near a weight endpoint, and the quadratic
metaplectic realization on the limiting oscillator space. Krawtchouk
polynomials become Hermite polynomials under central scaling; the
Ehrenfest generator becomes the Ornstein--Uhlenbeck generator, and
Gaussian conjugation gives the oscillator. The same construction
identifies finite bilinear kernels with precursors of Mehler kernels.
For convolution, a bulk--bulk limit loses its inputs, whereas a
bulk--macroscopic limit retains a nontrivial Ornstein--Uhlenbeck kernel.
We also identify the squeezed Gaussian kernel in the spectral working
note, distinguishing its exact diagonalization from the still conditional
transfer of its spectrum to finite matrices. The endpoint kernels
$b=\pm2a$ are treated at the end.
\end{abstract}
\maketitle
\tableofcontents

\section{Source map, scope, and principal distinctions}
This note correlates the following materials in the supplied Drive project.
\begin{description}
\item[FeinNotes11 (F).] The working transcription in
the FeinNotes11 research-note folder: Hermite wavefunctions,
Fourier phases, the quadratic $\mathfrak{sl}_2$ triple, and the question
of a convolution intertwiner.
\item[Encyclopedia (E).] Chapters 3 (urn chains), 6 (finite spin), 10
(Heisenberg algebra), 11 (central limits), 14 (transforms), 16 (symmetric
matrices), and 18 (association schemes). The Mehler subsection in
Chapter 11 is presently a placeholder; Sections~\ref{sec:finite-kernels}--\ref{sec:gaussian}
below supply a standalone development of that connection.
\item[Differential note (D).] The finite polynomial realization in
the Differential Realization research-note folder. Its differential
operators act on finite-degree auxiliary polynomials; they are not the
Schrodinger oscillator merely because both models use derivatives.
\item[Spectral note (S).] The working note in
the Spectral Working Note folder: the matrix $S_N=\Phi_NB_N$,
a commuting construction, local Gaussian asymptotics, and a candidate
squeezed Mehler limit.
\end{description}
These are internal source correlations, not new bibliographic claims.
F itself records that it is an editorial transcription, not a newly
verified diplomatic reading of the original scan.

The main logical chain is
\[
\begin{array}{ccccc}
\text{finite spin}&\longrightarrow&\text{Heisenberg algebra}
 &\longrightarrow&\text{quadratic oscillator algebra}\\
\text{Krawtchouk}&\longrightarrow&\text{Hermite}
 &\longrightarrow&\text{Hermite wavefunctions}\\
\text{Ehrenfest}&\longrightarrow&\text{Ornstein--Uhlenbeck}
 &\longrightarrow&\text{oscillator heat evolution}.
\end{array}
\]
The second arrow in the last row is Gaussian conjugation, not another
probabilistic scaling limit. The exact Mehler diagonalization of a
candidate limiting kernel is not, by itself, a proof of convergence of
finite matrix spectra.

\section{Hermite normalization and the Fourier sign audit}\label{sec:hermite}
Set $D=d/dx$ and define
\begin{equation}\label{eq:hermite-gen}
 e^{zx-\tau z^2/2}=\sum_{n\ge0}H_n(x,\tau)\frac{z^n}{n!}.
\end{equation}
We reserve $H_n(x)=H_n(x,1)$ for probabilists' Hermite polynomials.
In particular,
\[
 H_n(x,\tfrac12)=2^{-n/2}H_n(\sqrt2x).
\]
Let
\[
 \psi_n(x)=2^nH_n(x,\tfrac12)e^{-x^2/2},\qquad
 h_n(x)=\frac{\psi_n(x)}{\pi^{1/4}\sqrt{2^nn!}}.
\]
Their wavefunction generating function is
\begin{equation}\label{eq:wave-gen}
 \sum_{n\ge0}\psi_n(x)\frac{u^n}{n!}
 =e^{-x^2/2+2ux-u^2}.
\end{equation}
Differentiating \eqref{eq:hermite-gen} in $x$ and $z$ gives
\[
 H_n'=nH_{n-1},\qquad
 H_{n+1}(x,\tau)=xH_n(x,\tau)-n\tau H_{n-1}(x,\tau).
\]
Consequently,
\begin{equation}\label{eq:ladder}
 (x+D)\psi_n=2n\psi_{n-1},\quad
 (x-D)\psi_n=\psi_{n+1},\quad
 H\psi_n=(n+\tfrac12)\psi_n,
 \qquad H=\tfrac12(x^2-D^2).
\end{equation}
The normalized $h_n$ form an orthonormal basis of $L^2(\R,dx)$.
Write
\[
 a=\frac{x+D}{\sqrt2},\quad a^\dagger=\frac{x-D}{\sqrt2},\quad
 \NN=H-\tfrac12\Id=a^\dagger a.
\]
Then $[a,a^\dagger]=\Id$, $ah_n=\sqrt n\,h_{n-1}$, and
$a^\dagger h_n=\sqrt{n+1}\,h_{n+1}$.

\subsection{One convention used consistently}
Throughout this note,
\begin{equation}\label{eq:fourier-def}
 (\F f)(\xi)=\frac1{\sqrt{2\pi}}\int_\R e^{-ix\xi}f(x)\,dx,
 \qquad
 (\F^*f)(\xi)=\frac1{\sqrt{2\pi}}\int_\R e^{+ix\xi}f(x)\,dx.
\end{equation}
The Gaussian is fixed by both transforms. Applying the raising identity,
or integrating \eqref{eq:wave-gen}, gives
\[
 \F h_n=(-i)^nh_n,\qquad \F^*h_n=i^nh_n.
\]
Therefore
\begin{equation}\label{eq:fourier-osc}
 \boxed{\F=e^{-i\pi\NN/2},\qquad
 e^{i\pi H/2}=e^{i\pi/4}\F^*,\qquad
 e^{-i\pi H/2}=e^{-i\pi/4}\F.}
\end{equation}
The square root $\sqrt i$ means $e^{i\pi/4}$.
If the positive-exponent transform is instead named $\F$, the two names
must be exchanged everywhere. One must not retain the negative exponent
in \eqref{eq:fourier-def} while writing $e^{i\pi H/2}=\sqrt i\,\F$.
The current F transcription already has the correct negative-exponent
$\F$ and positive-exponent $\F^*$; the accompanying sign audit makes this
explicit rather than introducing a new sign reversal. Its corrected
Gaussian mixed term is $2i\xi u$, not $i\xi u$.

\section{Finite spin and the convention dictionary}\label{sec:spin}
The E and D conventions are
\[
 R_Ne_j=(N-j)e_{j+1},\quad L_Ne_j=je_{j-1},\quad
 \Lambda_Ne_j=(N-2j)e_j,
\]
with endpoint terms zero. Thus
\begin{equation}\label{eq:finite-brackets}
 [L_N,R_N]=\Lambda_N,\quad
 [\Lambda_N,R_N]=-2R_N,\quad [\Lambda_N,L_N]=2L_N.
\end{equation}
These are not literally F's signs
\[
 [\mathsf L,\mathsf R]=\rho,\qquad
 [\mathsf L,\rho]=2\mathsf L,\qquad [\rho,\mathsf R]=2\mathsf R.
\]
An exact dictionary is
\begin{equation}\label{eq:sign-dictionary}
 \mathsf L_N=-L_N,\qquad \mathsf R_N=R_N,\qquad \rho_N=-\Lambda_N.
\end{equation}
This distinguishes an algebra convention from an adjoint convention.

Put $B_N=\diag\binom Nj$. The coordinate metric in E is $B_N^{-1}$.
Conjugating by $B_N^{1/2}$ gives the orthonormal-coordinate matrices
\[
 \widehat R_N=B_N^{-1/2}R_NB_N^{1/2},\quad
 \widehat L_N=B_N^{-1/2}L_NB_N^{1/2},
\]
whose actions are
\[
 \widehat R_N|j\rangle=\sqrt{(N-j)(j+1)}|j+1\rangle,\quad
 \widehat L_N|j\rangle=\sqrt{j(N-j+1)}|j-1\rangle.
\]
Now $\widehat L_N=\widehat R_N^*$. Taking $\hbar=1$,
\[
 J_+=\widehat L_N,\quad J_-=\widehat R_N,\quad J_z=\Lambda_N/2
\]
is the spin-$N/2$ representation, with magnetic number $m=N/2-j$.
Its Casimir is $(N/2)(N/2+1)\Id$.

On auxiliary polynomials of degree at most $N$, the same unnormalized
finite model is
\[
 L_N=\partial_z,\qquad R_N=Nz-z^2\partial_z,\qquad
 \Lambda_N=N-2z\partial_z.
\]
The cancellation in $R_Nz^N=0$ is essential. This is D's finite
polynomial realization, not the oscillator acting on $L^2(\R)$.

\section{The contraction and the two oscillator algebras}\label{sec:contraction}
Define
\[
 a_N=\widehat L_N/\sqrt N,\quad
 a_N^\dagger=\widehat R_N/\sqrt N,\quad
 \NN_N=(N\Id-\Lambda_N)/2.
\]
Then
\begin{equation}\label{eq:contraction}
 [a_N,a_N^\dagger]=\Id-2\NN_N/N,
 \qquad [\NN_N,a_N]=-a_N,\quad
 [\NN_N,a_N^\dagger]=a_N^\dagger.
\end{equation}
For each fixed $j$, the two ladder coefficients tend to $\sqrt j$ and
$\sqrt{j+1}$. Under $|j\rangle\mapsto h_j$, this gives the oscillator
on the algebraic span of the Hermite basis. Convergence is on each
fixed finite set of levels, not uniformly over all $N+1$ levels.
Finite-dimensional commutators have trace zero, so an exact identity
$[a_N,a_N^\dagger]=\Id$ is impossible; \eqref{eq:contraction} shows
where the compensating high-level contribution resides.

\subsection{Why F's quadratic triple is a further construction}
On Schwartz space, F uses
\begin{equation}\label{eq:quadratic}
 \mathsf R=x^2/2,\qquad \mathsf L=D^2/2,\qquad \rho=xD+1/2,
 \qquad H=\mathsf R-\mathsf L.
\end{equation}
These satisfy F's brackets. They are quadratic expressions in the
limiting Heisenberg generators. Thus the logical route is a spin
contraction to Heisenberg, followed by formation of quadratic operators;
it is not a literal identification of F's triple with the finite spin
matrices before taking a limit.

Equivalently, put
\[
 K_+=\tfrac12(a^\dagger)^2,\quad K_-=\tfrac12a^2,\quad K_0=H/2.
\]
Then $[K_0,K_\pm]=\pm K_\pm$ and $[K_-,K_+]=2K_0$.
These operators preserve parity. The even and odd subspaces have lowest
$K_0$ weights $1/4$ and $3/4$, respectively. The quadratic oscillator
representation is therefore reducible into two infinite-dimensional
pieces, unlike one irreducible finite spin space. The complex Lie
algebra is the same type; the real forms, adjoints, and representations
must not be conflated.

\section{Krawtchouk polynomials and the finite transform}\label{sec:K}
Use ordinary coefficient normalization
\begin{equation}\label{eq:Kgen}
 (1+z)^{N-k}(1-z)^k=\sum_{n=0}^N K_n(k;N)z^n.
\end{equation}
Thus $K_0=1$, $K_1=N-2k$, and
\[
 K_2=\tfrac12\big((N-2k)^2-N\big).
\]
For $\pi_N(k)=2^{-N}\binom Nk$,
\begin{equation}\label{eq:Korth}
 \sum_{k=0}^N\pi_N(k)K_n(k;N)K_m(k;N)
 =\binom Nn\delta_{nm}.
\end{equation}
Set $p_{n,N}=K_n/\sqrt{\binom Nn}$. In E's matrix convention,
\[
 (\Phi_N)_{nk}=K_n(k;N),\quad \Phi_N^2=2^N\Id,\quad S_N=\Phi_NB_N.
\]
The genuinely orthogonal, symmetric transform is
\begin{equation}\label{eq:Q}
 Q_N=2^{-N/2}B_N^{-1/2}\Phi_NB_N^{1/2},\qquad Q_N^TQ_N=\Id,
\end{equation}
with entries, by weighted duality,
\[
 (Q_N)_{kn}=\sqrt{\pi_N(k)}\,p_{n,N}(k).
\]
\emph{$S_N$ is symmetric but not orthogonal.} Its spectral problem in S
is different from the finite oscillator diagonalized by $Q_N$.

\subsection{The Hermite limit}
Use the centered coordinate and mesh
\begin{equation}\label{eq:scale}
 x_k=\frac{N-2k}{\sqrt{2N}},\qquad \delta_N=\sqrt{2/N}.
\end{equation}
Putting $z=u/\sqrt{2N}$ in \eqref{eq:Kgen} gives, for $x_{k_N}\to x$,
\[
 \log\big((1+z)^{N-k_N}(1-z)^{k_N}\big)
 \longrightarrow xu-u^2/4.
\]
Hence for fixed $n$,
\begin{equation}\label{eq:KtoH}
 n!(2N)^{-n/2}K_n(k_N;N)\longrightarrow H_n(x,\tfrac12),
 \quad
 p_{n,N}(k_N)\longrightarrow\frac{H_n(\sqrt2x)}{\sqrt{n!}}.
\end{equation}
The local binomial limit is
\[
 \pi_N(k_N)/\delta_N\longrightarrow\pi^{-1/2}e^{-x^2}.
\]
Combining these yields the normalized wavefunction limit
\begin{equation}\label{eq:wave-limit}
 \delta_N^{-1/2}(Q_N)_{k_Nn}\longrightarrow h_n(x).
\end{equation}
If instead $x_k=(2k-N)/\sqrt{2N}$ is used, each degree-$n$ limit gains
$(-1)^n$. This is a coordinate reflection, not a new Fourier convention.

\section{Ehrenfest diffusion and Gaussian conjugation}\label{sec:diffusion}
The backward generator at fixed $N$ is
\begin{equation}\label{eq:GN}
 (G_Nf)(k)=\frac{N-k}{2}\big(f(k+1)-f(k)\big)
          +\frac k2\big(f(k-1)-f(k)\big).
\end{equation}
It is reversible for $\pi_N$, and $G_Np_{n,N}=-np_{n,N}$.
For a smooth function sampled in \eqref{eq:scale}, Taylor expansion gives
\[
 G_N\longrightarrow G=\tfrac12D^2-xD.
\]
Indeed the first-derivative coefficient is $-x$, and the second-derivative
coefficient is $1/2$. The invariant measure tends to
\[
 d\mu(x)=\pi^{-1/2}e^{-x^2}\,dx.
\]
This is the Ornstein--Uhlenbeck generator with stationary variance $1/2$.
The operator
\[
 U:L^2(d\mu)\to L^2(dx),\qquad
 (Uf)(x)=\pi^{-1/4}e^{-x^2/2}f(x)
\]
is unitary and satisfies
\begin{equation}\label{eq:OUconjugacy}
 -UGU^{-1}=H-\tfrac12\Id=\NN.
\end{equation}
At finite $N$, $U_Nf(k)=\sqrt{\pi_N(k)}f(k)$ gives
\[
 -U_NG_NU_N^{-1}
 =\frac N2\Id-\frac12(\widehat R_N+\widehat L_N)
 =:M_N,
\]
with $M_NQ_N=Q_N\diag(0,1,\ldots,N)$.
This is a rotated spin component shifted by $N/2$, not the unrotated
number matrix $\NN_N$ from Section~\ref{sec:contraction}; the two are
unitarily equivalent through $Q_N$.

The random-walk Hermite martingales in E, Chapter 11, describe a different
scaling: both the walk time and its polynomial order parameter change.
Their Brownian limit uses
$\partial_\tau H_n+\tfrac12\partial_x^2H_n=0$.
The fixed-$N$ urn chain above instead yields the stationary OU process.
Both are valid, but they are not the same finite dynamics.

\section{Finite kernels, Mehler kernels, and fractional transforms}\label{sec:finite-kernels}
For $s>0$ set $r=e^{-s}$. The finite Markov kernel relative to $\pi_N$ is
\begin{equation}\label{eq:finite-mehler}
 m_{r,N}(k,\ell)=\sum_{n=0}^N r^n p_{n,N}(k)p_{n,N}(\ell),
\end{equation}
so $P_s(k,\ell)=\pi_N(\ell)m_{r,N}(k,\ell)$.
It satisfies the exact finite composition identity
\[
 \sum_j\pi_N(j)m_{r,N}(k,j)m_{q,N}(j,\ell)
 =m_{rq,N}(k,\ell).
\]
This is the finite bilinear-kernel precursor missing from E's Mehler
placeholder.

At the process level, each of $N$ independent bits flips at rate $1/2$.
A bit has probability $(1-r)/2$ of changing after time $s$. Therefore
\[
 \sum_\ell P_s(k,\ell)z^\ell
 =\left(\frac{1+r}{2}+\frac{1-r}{2}z\right)^{N-k}
  \left(\frac{1-r}{2}+\frac{1+r}{2}z\right)^k.
\]
This formula gives a direct central-limit route to the OU transition
law, without interchanging an infinite spectral sum on the basis of
fixed-degree convergence alone.

\subsection{The limiting Markov and oscillator kernels}
In the limit,
\begin{equation}\label{eq:OUformula}
 (T_rf)(x)=\int_\R f(rx+\sqrt{1-r^2}\,z)\,d\mu(z).
\end{equation}
Its density relative to $\mu$ is
\begin{equation}\label{eq:OUkernel}
 m_r(x,y)=\frac1{\sqrt{1-r^2}}
 \exp\left(\frac{2rxy-r^2(x^2+y^2)}{1-r^2}\right).
\end{equation}
Conjugating gives the Lebesgue kernel of $r^{\NN}$:
\begin{equation}\label{eq:normalized-mehler}
 \sum_{n\ge0}r^nh_n(x)h_n(y)
 =\frac1{\sqrt{\pi(1-r^2)}}
 \exp\left[-\frac{1+r^2}{2(1-r^2)}(x^2+y^2)
            +\frac{2r}{1-r^2}xy\right].
\end{equation}
Thus the oscillator heat kernel is
\begin{equation}\label{eq:heat}
 (e^{-sH})(x,y)=\frac1{\sqrt{2\pi\sinh s}}
 \exp\left[-\frac{(x^2+y^2)\cosh s-2xy}{2\sinh s}\right].
\end{equation}
It is $e^{-s/2}$ times the kernel in \eqref{eq:normalized-mehler}.

For the requested probabilists' convention $H_n(x)=H_n(x,1)$,
\begin{equation}\label{eq:prob-mehler}
 \boxed{\sum_{n\ge0}\frac{r^nH_n(x)H_n(y)}{n!}
 =\frac1{\sqrt{1-r^2}}
 \exp\left(\frac{2rxy-r^2(x^2+y^2)}{2(1-r^2)}\right).}
\end{equation}
For real fixed $x,y$, the series converges absolutely for $|r|<1$,
locally uniformly in the disk; its radius is exactly $1$.
The square root equals $1$ at $r=0$. Boundary convergence is not included
in this assertion. Formula \eqref{eq:prob-mehler} uses the Gaussian
measure $(2\pi)^{-1/2}e^{-y^2/2}dy$ when read as a Markov kernel.
The factors of two in \eqref{eq:OUkernel} reflect the variance-$1/2$
coordinate used for our oscillator.

\subsection{Finite fractional transforms}
A phase-modified finite transform can be defined spectrally by
\[
 \mathfrak F_{N,\theta}=Q_N\diag(e^{-i\theta n})Q_N^T=e^{-i\theta M_N}.
\]
It is unitary, and on fixed limiting Hermite modes it tends to
$e^{-i\theta\NN}$. At $\theta=\pi/2$ this is $\F$.
This fixed-mode assertion does not claim pointwise convergence of an
oscillatory kernel everywhere. In particular $Q_N$ itself, which is
an involution, must not simply be identified with the continuum Fourier
transform, which has order four. D's cubic phase-modified composition
operator $\mathcal I\Phi_N$ is another finite group element, not by
notation the same matrix as $\mathfrak F_{N,\pi/2}$.

\section{Positive time, analytic continuation, and domains}
For $t<0$,
\begin{equation}\label{eq:t-kernel}
 (e^{tH})(x,y)=\frac1{\sqrt{-2\pi\sinh t}}
 \exp\left[\frac{(x^2+y^2)\cosh t-2xy}{2\sinh t}\right].
\end{equation}
It defines a bounded operator on all $L^2$, of norm $e^{t/2}$, and its
integral on each $\psi_n$ is absolutely convergent.
For $t>0$, the spectral operator is unbounded, with domain
\[
 \mathcal D(e^{tH})=
 \left\{\sum c_nh_n:\sum e^{2t(n+1/2)}|c_n|^2<\infty\right\}.
\]
It acts on each $h_n$ by $e^{t(n+1/2)}$. But substituting $t>0$ in
\eqref{eq:t-kernel} does not yield a convergent real-line integral,
even on those eigenfunctions: the resulting $y^2$ coefficient is
\[
 \tfrac12(\coth t-1)=\frac1{e^{2t}-1}>0.
\]
For $\operatorname{Re}t<0$, \eqref{eq:t-kernel} is analytic, with the
square root continued from the positive one on the negative real axis.
The boundary at $t=i\pi/2$ is
\[
 \frac{e^{i\pi/4}}{\sqrt{2\pi}}e^{ixy},
\]
recovering \eqref{eq:fourier-osc}. Such boundary kernels are interpreted
oscillatorily, or through strong limits, not as absolutely integrable
kernels acting on arbitrary $L^2$ functions.

\section{General Gaussian kernels and squeezing}\label{sec:gaussian}
Consider a kernel with a scalar amplitude $A$,
\[
 K(x,y)=A\exp[-a(x^2+y^2)+bxy],\qquad a>0,\quad |b|<2a.
\]
Set
\begin{equation}\label{eq:abmatch}
 \omega=\sqrt{4a^2-b^2},\qquad r=\frac{b}{2a+\omega},\qquad
 C=\sqrt{\frac{\pi(1-r^2)}{\omega}}.
\end{equation}
Then
\[
 a=\frac{\omega(1+r^2)}{2(1-r^2)},\qquad
 b=\frac{2\omega r}{1-r^2}.
\]
Define
\[
 H_\omega=\tfrac12(-D^2+\omega^2x^2),\quad
 \NN_\omega=H_\omega/\omega-\tfrac12\Id,\quad
 h_n^{(\omega)}(x)=\omega^{1/4}h_n(\sqrt\omega x).
\]
The exact expansion is
\begin{equation}\label{eq:gauss-spectral}
 K(x,y)=AC\sum_{n\ge0}r^nh_n^{(\omega)}(x)h_n^{(\omega)}(y).
\end{equation}
For $r>0$, this is $ACr^{\NN_\omega}$. For $r<0$, it is
$AC P|r|^{\NN_\omega}$, where $Pf(x)=f(-x)$. At $r=0$ it is the
rank-one ground-state projection times $AC$.
For real $A$, the eigenvalues are $ACr^n$; absolute summability makes
the operator trace class.

The dilation
\[
 (V_\omega f)(x)=\omega^{1/4}f(\sqrt\omega x)
\]
satisfies
\[
 \NN_\omega=V_\omega\NN V_\omega^{-1},\qquad
 V_\omega=e^{(\log\omega)\rho/2}.
\]
Here $\rho=xD+1/2$ is skew-adjoint after closure, so real dilation
parameters give unitary squeezing. Thus the Gaussian kernel is a
function of a squeezed oscillator, not generally a function of $H$ in
the original coordinates. No translation is needed for this centered
quadratic form.

For $0<b<2a$, write $r=e^{-\tau}$, $\tau>0$. Equivalently,
\[
 a=\frac\omega2\coth\tau,\qquad b=\frac\omega{\sinh\tau},
\]
and
\begin{equation}\label{eq:gauss-heat}
 K=A\sqrt{\frac{2\pi\sinh\tau}{\omega}}\,
 e^{-\tau H_\omega/\omega}.
\end{equation}

\subsection{A worked normalization check}
For $a=3/2$, $b=1$, one obtains
\[
 \omega=\sqrt8,\qquad r=3-\sqrt8,\qquad
 \tau=2\log(1+\sqrt2).
\]
The ground-state width is therefore $\sqrt8$, not the unsqueezed width
$1$. This is exactly the parameter combination occurring in S.

\section{Correlation with the spectral working note}\label{sec:spectral}
The candidate kernel in S, in its own centered coordinate
$y=(i-N/2)/\sqrt N$, is
\begin{equation}\label{eq:S-kernel}
 \mathcal G(y,z)=\frac{\sqrt2}{\pi^{3/2}}
 \exp[-\tfrac32(y^2+z^2)+yz].
\end{equation}
Its coordinate is related to \eqref{eq:scale} by $x=-\sqrt2y$;
frequency parameters must not be transferred without this rescaling.
Applying \eqref{eq:abmatch} to \eqref{eq:S-kernel} gives exactly
\begin{equation}\label{eq:S-identification}
 \boxed{\mathcal G=cq^{\NN_\alpha},\qquad
 \alpha=\sqrt8,\quad q=3-\sqrt8,\quad
 c=\frac{\sqrt8-2}{\pi}.}
\end{equation}
Hence the candidate integral operator has eigenfunctions
$h_n^{(\alpha)}$ and eigenvalues $cq^n$. Its Gaussian ground state is
an exact eigenfunction, not merely a variational approximation to
this integral operator. The equivalent Lie expression is
\[
 \mathcal G=c e^{\tau/2}
 \exp[\tau(\alpha^{-1}\mathsf L-\alpha\mathsf R)],
 \qquad \tau=-\log q.
\]
Thus F's algebra supplies the correct interpretation of S's squeezing
and imaginary-time evolution.

\subsection{What is proved, and what is not inferred}
S derives compact-window asymptotics on equal-parity indices of a
positive matrix associated with $S_N^2$. The spacing there is
$2/\sqrt N$, not $1/\sqrt N$; that factor enters the amplitude in
\eqref{eq:S-kernel}. It establishes a compactly supported Rayleigh
limit and the lower bound
\begin{equation}\label{eq:lower}
 \liminf_{N\to\infty}\frac{\pi N\lambda_{\max}(S_N)^2}{8^N}
 \ge\sqrt8-2.
\end{equation}
Equation \eqref{eq:S-identification} is an exact fact about the candidate
integral operator. It does not replace the missing global upper estimate
or prove that all finite eigenvalues converge to $cq^n$. Those claims
require control outside compact central windows, exclusion of competing
modes, and an appropriate operator or spectral convergence argument.
The ordinary finite-spin/OU limit and the weighted $S_N$ spectral
problem are related through oscillator structure, but are not the same
asymptotic theorem.

\section{Convolution: what survives, what collapses, and what must be specified}\label{sec:convolution}
F asks for an intertwiner between ordinary convolution/Fourier analysis
and Krawtchouk convolution/transform. There are several distinct questions
hidden here. An exact finite transform identity, a positive convolution
limit, and an abstract transport of a product are different assertions.

\subsection{Positive radial convolution on the Hamming cube}
Represent a cube point by a subset of $\{1,\ldots,N\}$ and add by symmetric
difference. If independent uniform subsets have sizes $k,\ell$, their
sum has size $M=k+\ell-2J$, where $J$ is the intersection size. Thus
\begin{equation}\label{eq:hypergroup}
 \delta_k*_N\delta_\ell
 =\sum_j\frac{\binom kj\binom{N-k}{\ell-j}}{\binom N\ell}
 \delta_{k+\ell-2j}.
\end{equation}
The sum uses precisely the feasible intersection indices. Its identity
is $\delta_0$, and its invariant probability is $\pi_N$.
The normalized characters are
\[
 \chi_{n,N}(k)=K_n(k;N)/\binom Nn.
\]
They obey the product formula
\begin{equation}\label{eq:character}
 \int\chi_{n,N}(m)\,d(\delta_k*_N\delta_\ell)(m)
 =\chi_{n,N}(k)\chi_{n,N}(\ell).
\end{equation}
For finite measures, $\widehat\nu(n)=\sum_k\chi_{n,N}(k)\nu(k)$ therefore
turns $*_N$ into pointwise multiplication. This is a precise finite
convolution theorem. It does not by itself identify that convolution
with ordinary addition on $\R$.

\subsection{Bulk--bulk scaling: loss of the two inputs}
Suppose both $x_{k_N}\to x$ and $x_{\ell_N}\to y$ in \eqref{eq:scale}.
The hypergeometric mean and variance imply
\[
 \mathbb E x_M=\sqrt{2/N}\,x_{k_N}x_{\ell_N}\longrightarrow0,
 \qquad \operatorname{Var}(x_M)\longrightarrow\tfrac12.
\]
The hypergeometric central limit theorem gives the pushforward limit
\begin{equation}\label{eq:collapse}
 \delta_{k_N}*_N\delta_{\ell_N}\Longrightarrow\mu.
\end{equation}
So two central inputs produce the invariant Gaussian, independently
of their limiting coordinates. There is no ordinary point-addition
convolution in this scaling. The finite identity $k=0$ has escaped to
$x=\sqrt{N/2}\to+\infty$.

The character calculation explains the same phenomenon: for fixed
$n>0$, $\chi_{n,N}(k_N)$ is of order $N^{-n/2}$. The renormalization
needed to obtain a nonzero Hermite polynomial does not preserve the
character product formula, because it occurs twice on the right side
and once on the left.

\subsection{Bulk--macroscopic scaling: a nontrivial Mehler limit}\label{sec:mixed}
There is an important refinement of the preceding collapse. Keep
$x_{k_N}\to x$ but take
\[
 \ell_N/N\longrightarrow(1-r)/2,\qquad -1<r<1.
\]
Then the same hypergeometric calculation gives
\[
 \mathbb E x_M\longrightarrow rx,\qquad
 \operatorname{Var}(x_M)\longrightarrow\frac{1-r^2}{2}.
\]
Its conditional central limit law is
\begin{equation}\label{eq:mixed}
 x_M\Longrightarrow rx+\sqrt{1-r^2}\,Z,\qquad Z\sim\mu.
\end{equation}
This is exactly \eqref{eq:OUformula}, including negative $r$ as a reflected
OU kernel. For $r=\pm1$, bounded central-scale variances collapse and the
corresponding transition limit is $x\mapsto\pm x$.

Spectrally, for every fixed $n$,
\[
 \chi_{n,N}(\ell_N)\longrightarrow r^n.
\]
Unlike the bulk--bulk case, one character factor now remains of order
one, so \eqref{eq:character} has a nontrivial Hermite limit.
Thus it is too strong to say that Krawtchouk convolution has no useful
continuum limit: it has a nontrivial action on the central variable
when the other input is kept at macroscopic scale.

There is even an exact finite connection. Convolution by the radial
probability
\[
 \nu_{r,N}(\ell)=\binom N\ell
 \left(\frac{1-r}{2}\right)^\ell
 \left(\frac{1+r}{2}\right)^{N-\ell}
\]
means independently flipping each coordinate with probability
$(1-r)/2$. Its character multipliers are $r^n$, so it is precisely the
finite kernel \eqref{eq:finite-mehler}. For $r=e^{-s}$, it is the
Ehrenfest semigroup. Moreover,
\[
 \nu_{r,N}*_N\nu_{q,N}=\nu_{rq,N}.
\]
The limiting parameter law is therefore multiplication of $r$'s, or
addition of the times $s=-\log r$ for positive parameters.

\subsection{Kernel convolution is preserved}
The limiting identity is
\begin{equation}\label{eq:kernel-compose}
 \int_\R m_r(x,z)m_q(z,y)\,d\mu(z)=m_{rq}(x,y).
\end{equation}
After Gaussian conjugation it is the semigroup law for $r^{\NN}$,
and after the scalar ground-energy correction it is
$e^{-sH}e^{-tH}=e^{-(s+t)H}$.
This is the most direct positive convolution connection between the
finite kernels and the oscillator.

\subsection{Ordinary convolution and Wick multiplication are different}
Ordinary convolution uses addition of independent real variables and
is diagonalized by $e^{ix\xi}$. OU evolution instead combines a
contracted input with independent Gaussian noise. The identity
\[
 H_n(ax+by,\tfrac12)
 =\sum_{j=0}^n\binom nj a^jb^{n-j}
 H_j(x,\tfrac12)H_{n-j}(y,\tfrac12),\qquad a^2+b^2=1,
\]
follows directly from the generating function. It expresses Gaussian
addition with degree splitting, not a same-degree character product.

On polynomials one may separately define Wick multiplication by
\[
 H_m(x,\tfrac12)\diamond H_n(x,\tfrac12)=H_{m+n}(x,\tfrac12).
\]
Writing $W=e^{-D^2/4}$ on polynomials, one has $Wx^n=H_n(x,\tfrac12)$
and
\[
 f\diamond g=W\big((W^{-1}f)(W^{-1}g)\big).
\]
This is a well-defined algebraic product because the derivative series
terminate on polynomials. It is not the positive radial convolution
\eqref{eq:hypergroup}, and no general function-space extension follows
without additional domain choices.

An arbitrary invertible finite transform also permits an abstractly
transported product. Such a construction alone says nothing about
positivity, locality, or a continuum limit. This qualifies the
intertwiner question in F: a natural finite convolution theorem exists,
and a natural OU/Mehler scaling exists, but an unqualified equivalence
to ordinary real-line convolution is not established by either one.

\section{A practical dictionary and a proof agenda}
\begin{center}\small
\begin{tabular}{p{0.28\textwidth}p{0.29\textwidth}p{0.30\textwidth}}
\toprule
Finite object & Limiting object & Qualification\\\midrule
Spin ladders $\widehat L_N,\widehat R_N$ divided by $\sqrt N$
& $a,a^\dagger$ & Fixed excitation levels\\
$K_n(k;N)$ with central normalization
& $H_n(x,1/2)$ & Fixed polynomial degree\\
$Q_N$ columns divided by $\sqrt{\delta_N}$
& $h_n(x)$ & Fixed column, central coordinates\\
Ehrenfest generator $G_N$ & $\tfrac12D^2-xD$ & Diffusion scaling\\
Symmetrized generator $M_N$ & $H-\tfrac12I$ & Gaussian conjugation\\
$r^n$ bilinear kernels & OU and Mehler kernels & Measure must be specified\\
Finite phases $e^{-i\theta n}$ & $e^{-i\theta\NN}$ & Fixed-mode contraction\\
Weighted matrix $S_N$ problem & Squeezed candidate $cq^{\NN_\alpha}$
& Spectral transfer requires more proof\\\bottomrule
\end{tabular}
\end{center}
For further development, the main tasks are: specify embeddings and
operator topologies for a full finite-transform convergence theorem;
prove global tightness or domination when going beyond compact kernel
windows; complete a matching upper bound in S; and state precisely
which convolution product, measure, and topology any proposed
intertwiner is to preserve. The mixed-scale observation in
Section~\ref{sec:mixed} provides a particularly concrete formulation of
the convolution question that leads directly to Mehler evolution.

\section{Endpoint kernels: $b=\pm2a$}\label{sec:endpoints}
We end with the boundary of the confining Gaussian region. For $a>0$,
let $A_{a,b}$ have the unit-amplitude kernel
$e^{-a(x^2+y^2)+bxy}$.

\subsection{The positive endpoint: free heat evolution}
At $b=2a$,
\[
 K(x,y)=e^{-a(x-y)^2}.
\]
The operator is ordinary convolution with a Gaussian. With
\eqref{eq:fourier-def},
\[
 \F(A_{a,2a}f)(\xi)=\sqrt{\pi/a}\,e^{-\xi^2/(4a)}\F f(\xi),
\]
so
\begin{equation}\label{eq:end-plus}
 \boxed{A_{a,2a}=\sqrt{\pi/a}\,e^{D^2/(4a)}.}
\end{equation}
Its norm is $\sqrt{\pi/a}$. It is bounded but noncompact on $L^2(\R)$;
it is neither Hilbert--Schmidt nor trace class. The Fourier multiplier
has continuous spectrum with closure $[0,\sqrt{\pi/a}]$, rather than a
geometric sequence of discrete Hermite eigenvalues.
Thus ordinary Gaussian convolution appears exactly when the confining
quadratic form degenerates.

\subsection{The negative endpoint: heat evolution and reflection}
At $b=-2a$,
\[
 K(x,y)=e^{-a(x+y)^2},\qquad
 \boxed{A_{a,-2a}=\sqrt{\pi/a}\,e^{D^2/(4a)}P.}
\]
Reflection commutes with the free heat operator. On even functions this
is the positive endpoint operator; on odd functions it is its negative.
It is bounded and noncompact with the same norm, and its spectrum is
$[-\sqrt{\pi/a},\sqrt{\pi/a}]$.

\subsection{How the squeezed oscillator reaches the endpoint}
As $b\uparrow2a$ with $a$ fixed,
\[
 \omega\to0,\quad r\to1,\quad \tau=-\log r\to0,
 \qquad \tau/\omega\to1/(2a).
\]
In \eqref{eq:gauss-heat}, the potential disappears and the kinetic term
approaches $D^2/(4a)$; the scalar tends to $\sqrt{\pi/a}$.
The kernel operators converge strongly to \eqref{eq:end-plus}. One way
to justify this is pointwise convergence on compactly supported inputs,
an integrable Gaussian tail bound for those inputs, and the uniform
Schur bound
\[
 \sup_x\int e^{-a(x^2+y^2)+bxy}\,dy
 =\sqrt{\pi/a},\qquad |b|\le2a.
\]
Density then gives strong convergence on $L^2$. The negative endpoint
follows by reflection. Convergence cannot be in operator norm: all
interior operators are compact, and both endpoint operators are not.
For $|b|>2a$, the quadratic form is indefinite and the displayed
integral kernel does not define a bounded operator on all of $L^2$.
\end{document}
